\documentclass[11pt,a4paper]{article}
\usepackage[margin=1in]{geometry}
\usepackage{titling}
\usepackage{xcolor}
\usepackage{listings}
\usepackage{hyperref}
\usepackage{longtable}
\usepackage{booktabs}
\usepackage{enumitem}
\usepackage{fancyhdr}
\usepackage{graphicx}

\definecolor{codebg}{rgb}{0.95,0.95,0.95}
\lstset{
  basicstyle=\ttfamily\footnotesize,
  backgroundcolor=\color{codebg},
  breaklines=true,
  frame=single,
  columns=fullflexible,
  showstringspaces=false
}

\pagestyle{fancy}
\fancyhf{}
\fancyhead[L]{MLOps AI Agent --- Technical Report}
\fancyhead[R]{\thepage}

\title{\textbf{Building a Production-Ready MLOps AI Agent}\\[0.3em]
\large Assignment No.\ 3 --- MLOps \& Production AI Engineering}
\author{Mohammad Hannan Abid \\ Roll No: 23i-0713 \\ Section: A}
\date{Fall 2026}

\begin{document}
\maketitle
\thispagestyle{empty}

\begin{center}
\textbf{Repository:} \url{https://github.com/HannanAbid7/assignment-03-mlops-agent}
\end{center}

\newpage
\tableofcontents
\newpage

\section{Introduction}

This report documents the design, implementation, and testing of an
AI-powered MLOps Agent capable of understanding natural-language requests
and operating four real MLOps systems: DVC (dataset versioning), MLflow
(experiment tracking), Apache Airflow (pipeline orchestration), and
Kubernetes (deployment management). The agent is built to answer only from
real tool evidence, never from its own training knowledge, and to never
silently trigger retraining or cluster-state changes.

\section{Complete System Architecture}

\begin{center}
\begin{verbatim}
User / Engineer (natural language)
        |
        v
FastAPI /chat endpoint (agent/main.py)
        |
        v
Orchestrator (agent/orchestrator.py)
  - sends system prompt + 21 tool schemas + conversation to an LLM
    (Groq, OpenAI-compatible API, model openai/gpt-oss-120b)
  - LLM decides which tool(s) to call
  - orchestrator executes the real tool function
  - result (status/data/error) fed back to the LLM
  - loop continues until a final answer (max 16 iterations,
    retries on HTTP 429 with real backoff)
        |
   -----+-----+-----+-----
   |    |     |     |
   v    v     v     v
  DVC  MLflow Airflow Kubernetes
\end{verbatim}
\end{center}

\subsection{Design decisions}

\begin{itemize}
  \item \textbf{No agent framework (e.g. LangChain) was used.} Tool calling
  is implemented directly against the OpenAI-compatible chat-completions
  API, so every tool invocation is explicit, inspectable in logs, and easy
  to reason about for the safety requirements below.
  \item \textbf{Tool schemas are auto-generated} from each Python function's
  real signature via \texttt{inspect.signature}, so the JSON schema sent to
  the LLM can never drift out of sync with the actual implementation
  (parameter names, types, and required/optional status are always
  consistent).
  \item \textbf{Every tool returns a uniform envelope}:
  \texttt{\{"status": "ok"|"error", "data": ..., "error": ...\}}, so the
  orchestrator can report failures honestly to the LLM instead of crashing
  or fabricating a result.
\end{itemize}

\section{Tool Definitions and Permissions}

Each of the four modules exposes 5--6 callable tools (21 total), satisfying
the "at least three tools per module" requirement. Table~\ref{tab:tools}
lists every tool with its real docstring (used verbatim as the tool
description sent to the LLM) and its safety/permission guard.

\begin{longtable}{p{2.1cm}p{4.3cm}p{7cm}}
\toprule
\textbf{Module} & \textbf{Tool} & \textbf{Permission / Safety Guard} \\
\midrule
\endhead
DVC & \texttt{list\_tracked\_datasets} & read-only \\
DVC & \texttt{get\_current\_version} & read-only \\
DVC & \texttt{compare\_versions} & read-only \\
DVC & \texttt{has\_data\_changed} & read-only \\
DVC & \texttt{get\_reproduction\_instructions} & read-only, no subprocess call \\
DVC (internal) & \texttt{\_run\_dvc} & subprocess restricted to an allow-list: \texttt{status}, \texttt{list}, \texttt{diff}, \texttt{checkout} only --- no destructive DVC command is reachable \\
\midrule
MLflow & \texttt{list\_experiments} & read-only \\
MLflow & \texttt{list\_runs} & read-only \\
MLflow & \texttt{get\_run\_details} & read-only \\
MLflow & \texttt{get\_best\_run} & read-only; always returns \texttt{production\_ready: false} with an explicit caveat so "highest metric" is never conflated with "safe to deploy" \\
MLflow & \texttt{compare\_runs} & read-only \\
MLflow & \texttt{get\_model\_versions} & read-only \\
\midrule
Airflow & \texttt{list\_dags} & read-only; filtered to this project's own DAGs (excludes $\sim$115 Airflow-bundled example DAGs) \\
Airflow & \texttt{get\_dag\_run\_status} & read-only \\
Airflow & \texttt{trigger\_dag} & \textbf{state-changing} --- allow-listed to \texttt{training\_pipeline} only; any other \texttt{dag\_id} raises \texttt{PermissionError} \\
Airflow & \texttt{get\_failed\_tasks} & read-only \\
Airflow & \texttt{get\_task\_logs} & read-only \\
\midrule
Kubernetes & \texttt{list\_pods\_and\_deployments} & read-only; restricted to \texttt{mlops-agent} namespace \\
Kubernetes & \texttt{get\_pod\_status\_and\_logs} & read-only; namespace-restricted \\
Kubernetes & \texttt{inspect\_services} & read-only; namespace-restricted \\
Kubernetes & \texttt{get\_unhealthy\_pods} & read-only; namespace-restricted \\
Kubernetes & \texttt{get\_deployed\_version} & read-only; namespace-restricted \\
Kubernetes & \texttt{restart\_deployment} & \textbf{state-changing} --- requires explicit \texttt{confirm=True} \emph{and} the approved namespace; silently refuses otherwise \\
\bottomrule
\caption{Complete tool inventory and permission model}
\label{tab:tools}
\end{longtable}

\subsection{System-level safety behaviors}

\begin{enumerate}
  \item The system prompt instructs the model to only state facts backed by
  a tool result, and to explicitly say when a tool errored or a service was
  unavailable, rather than guessing.
  \item Retraining and cluster-state changes are never triggered
  automatically: the system prompt instructs the model to only
  \emph{recommend} such actions, and the only two state-changing tools
  (\texttt{trigger\_dag}, \texttt{restart\_deployment}) carry their own
  independent guards regardless of what the model decides.
  \item Every tool call is logged with a request ID, tool name, status, and
  duration --- no secrets are ever logged.
  \item Malformed tool-call arguments from the LLM are caught and reported
  as a structured error instead of crashing the request.
\end{enumerate}

\section{Evidence from Each System}

\subsection{DVC}
\begin{lstlisting}
$ dvc status
Data and pipelines are up to date.

$ git tag -n
dataset-v1      fix: properly track dataset with DVC instead of git, ...
dataset-v2      chore: update dataset to v2
\end{lstlisting}
Two real dataset versions were created and tagged during development: training
on \texttt{dataset-v1} produced runs with F1 $\approx 1.0$; after removing
the last 10 rows of the Iris dataset to produce \texttt{dataset-v2}, F1
dropped to $\approx 0.93$ --- a genuine, observable effect of dataset
versioning on model quality, which the agent's \texttt{get\_current\_version}
and \texttt{compare\_versions} tools can detect.

\subsection{MLflow}
\begin{lstlisting}
run_id        params                              metrics
f4ab478858bb  n_estimators=100, max_depth=None     acc=0.929, f1=0.930
efe363d914a0  n_estimators=100, max_depth=None     acc=0.929, f1=0.930
2d3248b202ac  n_estimators=200, max_depth=10       acc=1.000, f1=1.000
dac909699dcf  n_estimators=150, max_depth=5        acc=1.000, f1=1.000
\end{lstlisting}
Eight manual training runs plus three Airflow-triggered runs were logged to
the \texttt{iris-training} experiment across both dataset versions, each
tagged with its producing git commit and DVC data hash.

\subsection{Airflow}
\begin{lstlisting}
$ GET /api/v2/dags/training_pipeline/dagRuns
manual__2026-11-11T18:08:43 -> success   (threshold 0.6, normal run)
manual__2026-11-11T18:49:54 -> success   (threshold 0.99, best historical
                                           run still cleared the gate)
manual__2026-11-11T18:50:31 -> failed    (threshold 1.5, quality gate
                                           correctly blocked promotion)
\end{lstlisting}
The DAG (\texttt{validate\_dataset} $\to$ \texttt{preprocess\_data} $\to$
\texttt{train\_model} $\to$ \texttt{evaluate\_and\_log\_mlflow} $\to$
\texttt{quality\_gate}) was triggered three times via the real Airflow REST
API (token-authenticated), demonstrating both a successful run and a
correctly-blocked run (Failure Injection 2, Section~\ref{sec:failures}).

\subsection{Kubernetes}
\begin{lstlisting}
$ kubectl get pods -n mlops-agent
NAME                        READY   STATUS    RESTARTS   AGE
iris-api-644d56d745-4kh6c   1/1     Running   0          157m
iris-api-644d56d745-k4tr2   1/1     Running   0          157m

$ kubectl get svc -n mlops-agent
NAME       TYPE        CLUSTER-IP      EXTERNAL-IP   PORT(S)
iris-api   ClusterIP   10.96.129.107   <none>        80/TCP

$ kubectl get deployment -n mlops-agent
NAME       READY   UP-TO-DATE   AVAILABLE   AGE
iris-api   2/2     2            2           157m
\end{lstlisting}
The inference service runs as a 2-replica Deployment with readiness/liveness
probes on \texttt{/health}, CPU/memory requests and limits, and
configuration via a ConfigMap, fronted by a ClusterIP Service --- all
required manifest elements are present (\texttt{kubernetes/*.yaml}).

\section{Tested Agent Conversations}

Three full conversations were run against the live agent (real LLM calls,
real tool execution against the running DVC/MLflow/Airflow/Kubernetes
stack). A fourth and fifth conversation were planned but could not be
completed with full final-answer text because the project's free-tier LLM
API key exhausted its daily token quota partway through testing --- this is
documented honestly as Known Limitation~1 (Section~\ref{sec:limitations})
rather than fabricated.

\subsection{Conversation 1: Mandatory multi-tool integration scenario}

\textbf{Prompt:} \textit{"Find the best recent model, identify the dataset
version used to train it, check the deployment status, and explain whether
retraining is necessary."}

\textbf{Tools called (in order):} \texttt{get\_best\_run} $\to$
\texttt{list\_pods\_and\_deployments} $\to$ \texttt{get\_deployed\_version}
$\to$ \texttt{compare\_versions} $\to$ \texttt{get\_model\_versions} $\to$
\texttt{list\_tracked\_datasets} $\to$ \texttt{get\_current\_version}.

\textbf{Agent's final answer (abridged):}
\begin{quote}
\small
\textbf{Known facts} --- Best recent run \texttt{2d3248b2...} (accuracy =
1.0, f1\_score = 1.0), tagged with \texttt{dataset\_version=v1} and DVC hash
\texttt{4d301ab...}. Current dataset hash \texttt{c49abde...} differs from
the run's hash --- the data has changed since that training. Deployment
\texttt{iris-api} is running version \texttt{iris-api:1.0} with 2 healthy
pods. Model registry is empty (no registered versions).

\textbf{Uncertainties} --- The deployed version has no explicit link to an
MLflow run ID; cannot confirm which model is actually serving.

\textbf{Recommendations} --- Register the validated model in the MLflow
Registry with its dataset hash and commit; since the current dataset differs
from the one used to train the best run, retraining on the latest data is
advisable --- but this is a \emph{recommendation only}, not an automatic
action.
\end{quote}

This demonstrates correlation of evidence across all four systems and
correctly stops short of auto-triggering retraining, per the assignment's
mandatory requirement.

\subsection{Conversation 2: Kubernetes deployment diagnosis (Failure Injection 1)}

\textbf{Prompt:} \textit{"Why is my prediction API unavailable?"}

\textbf{Injected failure:} \texttt{kubectl set image deployment/iris-api
iris-api=iris-api:bad -n mlops-agent}, producing a third pod stuck in
\texttt{ErrImagePull} alongside two healthy pods.

\textbf{Agent's final answer (abridged):}
\begin{quote}
\small
\textbf{Known facts} --- Two pods (\texttt{...4kh6c}, \texttt{...k4tr2}) are
\texttt{Running} with logs showing successful \texttt{GET /health} calls. A
third pod (\texttt{...5pnzc}) is \texttt{Pending}: "trying and failing to
pull image" --- a container-image-pull failure.

\textbf{Recommendations} --- Inspect the failing pod's events
(\texttt{kubectl describe pod}), validate the image reference in the
Deployment spec, fix the tag/registry issue, and roll out again; consider
temporarily scaling to the two healthy replicas while the image is fixed.
\end{quote}

The agent correctly distinguished healthy from unhealthy pods using real
pod status and logs, rather than guessing at the cause.

\subsection{Conversation 3: Airflow pipeline diagnosis (Failure Injection 2, partial)}

\textbf{Prompt:} \textit{"The last training pipeline run seems to have
failed, what happened and why?"}

\textbf{Injected failure:} \texttt{QUALITY\_THRESHOLD} raised from 0.6 to
1.5 (unreachable), DAG triggered.

\textbf{Tool chain executed (verified in logs):} \texttt{list\_dags} $\to$
\texttt{get\_dag\_run\_status} $\to$ \texttt{get\_failed\_tasks} $\to$
\texttt{get\_task\_logs}. The agent correctly identified that all four
upstream tasks succeeded and \texttt{quality\_gate} was the failed task,
fetching its logs as evidence.

\textbf{Outcome:} the LLM's final natural-language synthesis was cut off
mid-generation because the project's Groq API key hit its 200{,}000
token/day cap during this exact test (see Limitation~\ref{sec:limitations}).
The tool-selection and evidence-gathering logic itself is verified correct
from the structured logs; only the very last synthesis step did not
complete. This is reported honestly rather than fabricated, and is itself
documented as Failure Investigation 2.

\section{Failure Investigations}
\label{sec:failures}

\subsection{Failure Investigation 1: Kubernetes bad image tag}
\textbf{Root cause:} Deployment updated to reference a non-existent image
tag (\texttt{iris-api:bad}).\\
\textbf{Resolution:} \texttt{kubectl rollout undo deployment/iris-api -n
mlops-agent}; pods returned to \texttt{2/2 Running} within seconds.

\subsection{Failure Investigation 2: Airflow quality gate (plus a real external-API failure)}
\textbf{Root cause (intended):} \texttt{QUALITY\_THRESHOLD} set above any
achievable F1 score, so the quality gate correctly blocked promotion of an
otherwise-successfully-trained model.\\
\textbf{Root cause (unintended, discovered during this investigation):} the
Groq free-tier API key's \emph{daily} token quota (200{,}000 TPD) was
exhausted mid-test, which is a genuine infrastructure limitation of the
chosen LLM provider, separate from the agent's own logic.\\
\textbf{Resolution:} \texttt{QUALITY\_THRESHOLD} reverted to \texttt{0.6};
the daily-quota issue is documented as a known limitation (upgrading to a
paid tier or adding proactive token-budget tracking are the proposed fixes).

\section{Testing, Reliability, and Safety}

Ten automated \texttt{pytest} tests cover all four tool modules plus the
orchestrator, with every external service mocked:

\begin{itemize}
  \item \texttt{test\_rejects\_non\_allowlisted\_subcommand} --- DVC command allow-list enforced.
  \item \texttt{test\_get\_best\_run\_picks\_highest\_metric}, \texttt{test\_get\_best\_run\_no\_matches} --- MLflow tool correctness.
  \item \texttt{test\_trigger\_dag\_rejects\_unapproved} --- Airflow DAG allow-list enforced.
  \item \texttt{test\_get\_dag\_run\_status\_timeout} --- Airflow tool handles a connection timeout without crashing.
  \item \texttt{test\_rejects\_wrong\_namespace} --- Kubernetes namespace restriction enforced.
  \item \texttt{test\_restart\_requires\_confirm} --- unauthorized cluster-state change rejected.
  \item \texttt{test\_no\_invented\_results\_on\_tool\_failure} --- orchestrator reports a tool failure as an error, never a fabricated result.
  \item \texttt{test\_tool\_success\_path} --- successful tool execution and result propagation.
  \item \texttt{test\_logs\_contain\_no\_secrets} --- structured logs never expose credentials.
\end{itemize}

All 10 tests pass (\texttt{pytest -v}, verified locally and in CI via
GitHub Actions on every push and pull request).

Two deliberately injected failures were demonstrated end-to-end
(Section~\ref{sec:failures}), satisfying the requirement to show at least
two and explain how the agent diagnoses them.

\section{GitHub Repository and CI}

The repository follows the assignment's recommended structure
(\texttt{agent/}, \texttt{api/}, \texttt{dags/}, \texttt{kubernetes/},
\texttt{tests/}, \texttt{.github/workflows/}). A GitHub Actions workflow
(\texttt{.github/workflows/ci.yml}) runs linting and the full pytest suite
on every push and pull request, with all external service calls mocked so
CI requires no live infrastructure.

\section{Known Limitations and Proposed Improvements}
\label{sec:limitations}

\begin{enumerate}
  \item \textbf{Third-party API rate limits.} The free Groq tier enforces
  an 8{,}000 token/minute and 200{,}000 token/day cap. Long multi-tool
  conversations (especially ones pulling Airflow task logs) can exceed
  this, causing slow responses while the orchestrator backs off and
  retries, or --- once the \emph{daily} cap is hit --- a hard failure that
  cannot be retried around (observed directly in Conversation 3 / Failure
  Investigation 2). \textbf{Improvement:} move to a paid tier or
  self-hosted model, or add a proactive token-budget tracker.
  \item \textbf{Tool-result truncation can hide information.} Tool output
  is capped at 1{,}500 characters before being sent back to the LLM to
  control token usage; this can cut off relevant content in very large
  responses. \textbf{Improvement:} summarize large payloads server-side
  instead of blind truncation.
  \item \textbf{The training DAG's quality gate evaluates the best
  historical MLflow run, not strictly the run just produced}, since runs
  are not filtered by the triggering Airflow run ID. This means the gate
  can pass even if the \emph{current} training run underperforms a past
  one. \textbf{Improvement:} tag each MLflow run with its Airflow
  \texttt{run\_id} and filter the quality gate on that specific run.
  \item \textbf{No formal model registry workflow.} Models are logged to
  MLflow runs but never registered/promoted automatically --- this is
  intentional (the assignment requires that promotion never happen
  automatically), but it means the agent cannot yet answer "what model is
  in production" from the registry alone.
  \item \textbf{Single-node demo topology.} Airflow runs via
  \texttt{airflow standalone} and Kubernetes runs on a local \texttt{kind}
  cluster --- appropriate for this assignment, not representative of a
  production deployment.
\end{enumerate}

\section{Conclusion}

The MLOps AI Agent successfully integrates DVC, MLflow, Airflow, and
Kubernetes behind a single conversational interface, answers only from real
tool evidence, correctly withholds automatic retraining/cluster changes
behind explicit safety guards, and was validated against two deliberately
injected failures plus a mandatory multi-tool integration scenario. The one
incomplete test run was caused by an external API quota limit rather than a
flaw in the agent's own reasoning or tool-selection logic, and is documented
transparently above as both a limitation and a real-world engineering
lesson about depending on third-party rate-limited services.

\end{document}
